% Blocco sul tavolo legato, attraverso una carrucola sul bordo, a un pesetto
% appeso (problema ud2-dinamica-tavolo). Sul blocco: la tensione F_T in avanti
% e l'attrito dinamico F_d all'indietro (peso e forza normale si bilanciano e
% non sono disegnati). Sul pesetto: la tensione F_T verso l'alto e il peso
% m_2 g verso il basso. Le frecce a indicano il verso del moto, che si prende
% come positivo per ciascun corpo. Le forze sono in scala, 1 cm ogni 10 N:
% F_T = 9,6 N, m_2 g = 10 N, F_d = 8,0 N.
\begin{tikzpicture}[>=Stealth, font=\small]
  % tavolo
  \fill[richiamo!20] (-0.5,-0.6) rectangle (5.0,0);
  \draw[thick] (-0.5,0) -- (5.0,0) -- (5.0,-0.6);
  % carrucola sul bordo
  \coordinate (P) at (5.25,0.55);
  \draw[thick] (5.0,0) -- (P);
  \filldraw[fill=richiamo!25, draw=black, thick] (P) circle (0.3);
  \fill (P) circle (1.2pt);
  % blocco
  \filldraw[fill=esempio!20, draw=esempio, thick] (1.0,0) rectangle (2.4,0.9);
  \node[esempio] at (1.7,0.45) {$m_1$};
  % fune
  \draw[thick] (2.4,0.85) -- (5.25,0.85);
  \draw[thick] (5.55,0.55) -- (5.55,-1.4);
  % pesetto
  \filldraw[fill=esempio!20, draw=esempio, thick] (5.25,-1.4) rectangle (5.85,-2.0);
  \node[esempio] at (5.55,-1.7) {$m_2$};
  % forze sul blocco
  \draw[->, formula, very thick] (2.4,0.45) -- (3.36,0.45) node[right] {$\vec F_T$};
  \draw[->, medica, very thick] (1.0,0.45) -- (0.2,0.45) node[above] {$\vec F_d$};
  % forze sul pesetto
  \draw[->, formula, very thick] (5.85,-1.45) -- (5.85,-0.49) node[right] {$\vec F_T$};
  \draw[->, attenzione, very thick] (5.55,-2.0) -- (5.55,-3.0) node[right] {$m_2\vec g$};
  % versi del moto
  \draw[->, black!60, thick] (1.0,1.25) -- (2.4,1.25) node[right] {$a$};
  \draw[->, black!60, thick] (4.85,-1.4) -- (4.85,-2.4) node[left] {$a$};
\end{tikzpicture}
